\subsection{Standard Topology on Minkowski Space}\label{subsec:standard-topology-on-minkowski-space}
Physicists usually think of Minkowski space as the set $\mathbb{R}^4$ equipped with a \href{https://ncatlab.org/nlab/show/Euclidean+topology}{Euclidean topology}, a \href{https://ncatlab.org/nlab/show/Minkowski+metric}{Minkowski metric}, and the ``standard'' time-orientation. The ``standard'' time-orientation in this context is simply the non-vanishing vector $(1,0,0,0)$ in the standard coordinate system on $\mathbb{R}^4$.

This may be what Haag and Kastler mean by ``Minkowski space'', but then again it may not. As they don't explicitly define the term ``Minkowski space'', it's uncertain what they mean. Let's consider some alternatives.
